%% 第24回編集委員会（物理第3巻）：事務局
%% 物理シリーズ　第3巻　波動　巻末（こたえ・記号カード・音速・屈折率・物性値・なぜ？の道すじ・チェック表・つぎの巻へ・さくいん・奥付）
%% 第2巻の ph2-back.tex と同じ部品・同じ形。見出し（\bkbacktitle）の名前は、もくじ（ph3-toc.tex）の巻末の項目の名前と同じにする（tools/check_pages.py p3 が照らし合わせる）。
%% 第2巻の確認で決めたこと（第23回 U-1〜U-13、Q6）を、はじめから当てはめた：
%%   ・記号カード：数学の記号の行は、この巻の章の \bkkigocard（自動でうしろに集まる）と同じ記号を置かない（U-6・Q6-10）。
%%   ・紙面に委員会の回を書かない（U-7・Q6-11）。回はこのコメントにだけ書く。
%%   ・さくいんの段は \raggedcolumns。表4 の「先に読む本」は共通部品（bk-phys.sty）が巻の番号から自動で作るので、奥付をその文と同じにした（Q6-22）。
%% 物性値の出典は、各章のファイルの頭のコメント（出典の欄）と docs/作業メモ/物理第3巻_組どうしの確かめ.md の「巻末の表」の行から集め、
%%   事務局が次のものを作るときに確かめ直した（手引き 2章の13。11月5日）。計算は checks/物理第3巻/巻末/calc_ph3_back.py（出力 calc_ph3_back_out.txt）。
%%   ・光速・標準重力加速度・標準大気圧・ボルツマン定数：NIST の CODATA の値のページ（physics.nist.gov/cgi-bin/cuu/Value?c など）を curl で読んだ。
%%   ・水（8 ℃・20 ℃）とヘリウム（20 ℃）の密度と音速：NIST Chemistry WebBook（流体の熱物性、0.101325 MPa）を curl で読んだ。空気は CoolProp 8.0。
%%   ・屈折率：一次資料の分散式（水 Daimon・Masumura, Appl. Opt. 46, 3811 (2007)、ダイヤモンド Peter, Z. Phys. 15, 358 (1923)、
%%       N-BK7 SCHOTT のカタログ、MgF2 Dodge, Appl. Opt. 23, 1980 (1984)、空気 Ciddor, Appl. Opt. 35, 1566 (1996)）の係数から計算した。
%%       係数は refractiveindex.info のデータベース（GitHub の写し）から読んだ。水の 656 nm は、Daimon・Masumura の式で 1.3315（本文 1.331。Tilton・Taylor を引く表の値）。
%%   ・ナトリウムの D 線 589.0・589.6 nm、ネオンの 632.8 nm（ヘリウムネオンレーザー）：NIST Handbook of Basic Atomic Spectroscopic Data の表を curl で読んだ。
%%   ・A 特性：IEC 61672-1:2013 の附属書の式で計算し、表3 を写した表（python-acoustics の data/iec_61672_1_2013.csv、curl で読んだ）と小数第1位まで一致。
%%   ・騒音に係る環境基準：環境省 www.env.go.jp/kijun/oto1-1.html を curl で読んだ（AA 50/40、A・B 55/45、C 60/50、幹線道路に近接する空間 70/65 dB）。
%%   ・騒音障害防止のためのガイドライン（85 dB）：厚生労働省のパンフレット www.mhlw.go.jp/content/001195349.pdf を curl と pdftotext で読んだ。
%%   ・板ガラスの E 72 GPa・ν 0.23・2500 kg/m³：Pilkington の技術資料 ATS-129（WebSearch で、ほかの会社のデータシートと照らし合わせた）。
%%   ・ほかの値は、各章の組が確かめた出典（章の頭のコメント）のまま。Wikipedia は出典にしていない。
%%   〔要確認〕は、一次資料でも、一次資料を引く2つ以上の資料でも確かめられなかった値（家庭の洗濯機の脱水の回転数だけ。波-2）。
%% 表の書式（見出しの行の色・横の線・行の間）は、各表の colspec のあとにじかに書いた。第2巻の \ph@tblstyle のように
%%   マクロにまとめると、tabularray が展開せず、書式が効かない（試し組みで確かめた）。
\makeatletter
\def\phbkhead#1{\par\addvspace{4mm}\Needspace{6\baselineskip}\noindent{\sffamily\bfseries\large\color{bkDeep}#1}\par\nobreak\vspace{2mm}}
\newcommand\phyk{〔要確認〕}
%% ---------------- こたえ ----------------
\bkAnswerPages
\par\vspace{4mm}
%% 箱はページをまたがせない（またぐと、次のページに箱の切れはしだけが残る）
\begin{tcolorbox}[bk base, unbreakable, colframe=bkTheme!60, colback=bkTheme!5, boxrule=0.8pt, arc=2mm, left=4mm, right=4mm, top=2mm, bottom=2mm]
{\sffamily\bfseries\color{bkDeep}答え合わせのしかた}\par\vspace{1mm}
\begin{enumerate}[label={\bkcircled{\arabic*}}, leftmargin=2.4\zw, itemsep=0pt, topsep=0pt]
\item まず自分の答えを全部書いてから、ここを見よう。
\item 単位と桁を確かめよう。$f$〔Hz〕と $\omega=2\pi f$〔rad/s〕、振幅と強さ（振幅の2乗に比例）を取りちがえていないかも。
\item デシベルは、強さ・パワーの比なら $10\log_{10}$、音圧・振幅の比なら $20\log_{10}$。丸め方で最後の桁が1つずれても、間違いではありません。
\item ちがったら、同じ形の例題にもどって、どの段でちがったかを探そう。$c$ が減衰係数・波の速さ・光速のどれかも確かめよう。
\end{enumerate}
\end{tcolorbox}

%% ---------------- 記号カード ----------------
%% 数学シリーズの記号で、この巻でよく使うもの（この巻の章の \bkkigocard は自動でうしろに集まる）。
%% 章のカードにある記号（T・f、ω_n・A、c、ζ・δ、X、Q、L、q、M、K、μ、η、K_p、T_i・T_d、ε、τ、Z、I、r、v_s・v_o、a、L_p、t_R、d_33、H、Δf、b、C、n、f、E、m、θ_B）は、ここに置かない（U-6）。
\bkKigoPages{%
\bkkigorow{$\Delta$}{デルタ}{変化量・きざみ。時間のきざみ $\Delta t$、経路の差、温度差 $\Delta T$〔K〕}{2-3}%
\bkkigorow{rad}{ラジアン}{角の単位。$\pi$\,rad＝180°。位相のずれも rad で書く}{4-2}%
\bkkigorow{$\omega$}{オメガ}{角速度・角振動数〔rad/s〕。$\omega=2\pi f$}{4-8}%
\bkkigorow{$\phi$}{ファイ}{初期位相・位相のずれ。$A\sin(\omega t+\phi)$}{5-15}%
\bkkigorow{$\dot x$、\ $\ddot x$}{エックス・ドット、エックス・ツードット}{時間で1回微分（速度）・2回微分（加速度）。$m\ddot x+c\dot x+kx=F$}{5-5}%
\bkkigorow{$\approx$}{ニアリーイコール}{ほぼ等しい（近似）。ふれが小さいとき $\sin\theta\approx\theta$（5-17）}{5-8}%
\bkkigorow{$\displaystyle\int$、\ $\displaystyle\sum$}{インテグラル、シグマ}{細かく分けてたす・全部たす。フーリエ級数、DFT}{5-10}%
\bkkigorow{$e$、\ $\ln x$}{イー、エル・エヌ・エックス}{ネイピア数と自然対数。減っていくゆれ $e^{-\zeta\omega_{\mathrm n}t}$}{5-14}%
\bkkigorow{$\log_{10}x$、\ dB}{ログ、デシベル}{常用対数と、比を「10×常用対数」で表す単位。$\log_2$ はビットの数（波-11）}{3-12}%
\bkkigorow{$P(A)$}{ピー・エー}{起こる確率。情報量 $-\log_2 p_i$（波-11 は小文字の $p_i$）}{3-17}%
\bkkigorow{$j$}{ジェイ}{虚数単位（2乗すると $-1$）。高校・数学の本の $i$。電流の $i$ と区別する}{6-8}%
\bkkigorow{$\arg z$}{アーギュメント}{偏角。複素振幅の角が、ゆれのおくれ（負の角）}{6-9}%
\bkkigorow{$e^{j\theta}$}{イー・ジェイ・シータ}{単位円の上の角 $\theta$ の点（オイラーの公式）。回る矢印でゆれを書く}{6-10}%
\bkkigorow{$\det$}{デターミナント}{行列式。固有振動数を決める $\det(\mathbf K-\omega^2\mathbf M)=0$（波-3）}{7-4}%
\bkkigorow{$\lambda$、\ $\vec v$}{ラムダ、ベクトル・ブイ}{固有値と固有ベクトル。波-3 では $\lambda=\omega^2$ とモードの形（波長の $\lambda$ とは別）}{7-8}%
\bkkigorow{$\tfrac{\partial y}{\partial t}$}{ラウンドディー}{偏微分。波動方程式 $\tfrac{\partial^2y}{\partial t^2}=c^2\tfrac{\partial^2y}{\partial x^2}$}{8-2}%
\bkkigorow{$\nabla^2$}{ラプラシアン}{3つの向きの2階の偏微分の和。球面波・膜の振動}{8-12}%
\bkkigorow{$a_n$・$b_n$}{エー・エヌ、ビー・エヌ}{フーリエ係数。弦の音にふくまれる倍音の大きさ}{9-1}%
\bkkigorow{$F(\omega)$、\ $f_{\mathrm s}$}{エフ・オメガ、エフ・エス}{フーリエ変換（スペクトル）と、サンプリング周波数〔Hz〕}{9-4}%
\bkkigorow{$\mathcal L$、\ $s$}{ラプラス変換、エス}{ラプラス変換。微分が「$s$ をかける」になる（波-4）}{9-5}%
\bkkigorow{$\delta(t)$}{デルタ関数}{面積1で幅が限りなく細いパルス。たたいた力積（波-10 のインパルスハンマー）。対数減衰率 $\delta$ とは別}{9-6}%
\bkkigorow{$G(s)$}{ジー・エス}{伝達関数。出てくるもの÷入れるもの（はじめの値0）。$s=j\omega$ で周波数応答}{9-8}%
}

\Needspace{18\baselineskip}\phbkhead{分野で記号がぶつかるもの}
{\small 波動は、振動・波・音・光と分野が広く、同じ文字がちがう量を表すことがあります。この本では、同じ章の中では1つの文字を1つの意味にだけ使い、ぶつかるときは文字や添字を変えました。どの量かは、章（部）と単位で見分けよう。右から2つめの列は、第1巻『力学』・第2巻『熱学』とほかの巻での意味です。\par}\vspace{2mm}
{\footnotesize
\begin{longtblr}{colspec={Q[c,m,wd=15mm] X[l,m] Q[l,m,wd=34mm] Q[l,m,wd=28mm]}, width=\linewidth, rowhead=1, row{1}={bg=bkVol!14, font=\sffamily\bfseries\footnotesize}, hlines={0.4pt,bkLine}, rowsep=1.4pt, colsep=3.5pt}
記号 & この巻での意味 & 第1・2巻・ほかの巻での意味 & 見分け方 \\
$c$ & 減衰係数〔N·s/m〕（第1部：波-1〜4）、波の速さ・音速〔m/s〕（第2・3部：波-5〜11）、真空の光速〔m/s〕（第4部：波-12〜14）。波-7 ではダンパーを $b$、波-10 のセンサーは $\zeta$・$\omega_{\mathrm n}$ で書く & ダンパーの係数（力-2）、比熱（熱-1）。熱-13 は光速を $c_0$。熱-4 は音速を $a$ & 部で見分ける。単位：N·s/m、m/s \\
$k$ & ばね定数〔N/m〕（第1部）、波数 $2\pi/\lambda$〔rad/m〕（第2部から）。同じ章に両方出るときは、ばね定数を $k_{\mathrm s}$（波-5 ⑦・波-7）。波-13 は波数を使わず $2\pi/\lambda$ & ばね定数（力-2）、熱伝導率（熱-11）。ボルツマン定数は $k_{\mathrm B}$ & 単位：N/m、rad/m \\
$I$ & 慣性モーメント〔kg·m$^2$〕（波-1・3）、断面二次モーメント（波-6 ⑨）、波・音・光の強さ〔W/m$^2$〕（波-7・9・13）、情報量〔ビット〕（波-11）、光度〔cd〕（波-12 ⑩）。電流は記号で書かない（波-7） & 慣性モーメント（力-6）、断面二次モーメント（力-10）、電流（電-4）、単位行列（7-3） & 単位：kg·m$^2$、W/m$^2$、cd \\
$Z$ & インピーダンス（押すもの÷流れるもの。波-7・10）：弦 N·s/m、音 Pa·s/m（音響インピーダンス $\rho c$）、電気 Ω & 断面係数（力-10）、圧縮係数・分配関数（熱-3・熱-15）、電気のインピーダンス（電-10） & 単位で見分ける \\
$n$ & 屈折率 $c/v$（第4部）、モードの番号（波-6）、軸受けの玉の数（波-10） & 物質量（熱-3） & 屈折率は単位なし。第4部だけ \\
$T$ & 周期〔s〕（波-2・4・9・11 は使わず、$f$・$\omega$ で書く）、絶対温度〔K〕（波-9 だけ）、PID の積分時間・微分時間 $T_{\mathrm i}$・$T_{\mathrm d}$（波-4）。張力は $F_{\mathrm T}$（波-1・5・6・7） & 周期・ねじりモーメント（力-1・力-11）、絶対温度（熱-1） & 単位：s、K \\
$L$ & ラグランジアン $K-U$〔J〕（波-3）、むだ時間〔s〕（波-4）、距離・長さ（波-10・13）、インダクタンス〔H〕（相似の表・波-2）。レベルは $L_p$・$L_I$・$L_W$〔dB〕（波-9） & 長さ（力-8）、潜熱（熱-1）、インダクタンス（電-8） & 単位：J、s、m、H、dB \\
$\tau$ & 時定数〔s〕（波-4・11）、振幅の透過係数（波-7）、強さの透過率（波-9）、光学的な厚さ（波-13）。モーターのトルクは $\tau_{\mathrm m}$（波-4） & せん断応力（力-8）、時定数・透過率（熱-11・熱-13） & 時定数は s、ほかは単位なし \\
$H$ & ゆれの倍率（波-10）、情報のエントロピー〔ビット〕（波-11）、すきまの高さ〔m〕（波-13） & エンタルピー（熱-4）、全ヘッド（力-13）、ステップ関数 $H(t)$（9-5） & 単位：なし、ビット、m \\
$\mu$ & 質量比（波-3）、線密度〔kg/m〕（波-5〜7） & 摩擦係数・粘度（力-2・力-14）、化学ポテンシャル（熱-10） & 質量比は単位なし \\
$f$ & 振動数〔Hz〕（全章）、焦点距離〔m〕（波-12。振動数を $f$ と書かない） & 関数の名前（数学） & 単位：Hz、m \\
$E$ & ヤング率〔Pa〕（波-6・9）、照度〔lx〕（波-12）、光の電場の振幅 $E_0$（波-13）、1\,m あたりのエネルギー $\overline E$（波-7） & ヤング率（力-8）、電場（電-1） & 単位：Pa、lx \\
$m$ & 質量、干渉の次数（波-8・13。電子の質量は $m_{\mathrm e}$）、壁の面密度〔kg/m$^2$〕（波-9） & 質量（力-2） & 次数は整数 \\
$A$ & 振幅（全章。複素振幅は $X$）、吸音力〔m$^2$〕（波-9）、光線の行列の成分 $A$〜$D$（波-12）、断面積 & 断面積（力-8）、振幅は数学 6-6 では $X$ & 吸音力は m$^2$ \\
$K$、$\mathbf K$ & 運動エネルギー（波-1・3）、こわさ行列 $\mathbf K$（立てた太字。波-3）、比例ゲイン $K_{\mathrm p}$（波-4）、体積弾性率〔Pa〕（波-5・9）、平均する回数（波-11） & 運動エネルギー（力-4）、体積弾性率（力-8） & 太字・添字・単位 \\
$M$、$\mathbf M$ & 質量行列 $\mathbf M$（立てた太字。波-3）。台車の質量は細い字の $m$ & モル質量 $M$（熱-3） & 太字かどうか \\
$Q$ & Q 値（共振の鋭さ。波-2）、一般化力（波-3）、電荷（波-10）、体積流量（相似の表） & 熱量（熱-1）、体積流量（力-13） & Q 値は単位なし \\
$C$ & 静電容量〔F〕（相似の表・波-2）、部屋の熱容量〔J/K〕（波-4）、通信路容量〔ビット/s〕（波-11） & 熱容量 $C_{\mathrm{th}}$（熱-1）、静電容量（電-3） & 単位：F、J/K、ビット/s \\
$R$ & 抵抗〔Ω〕（相似の表）、熱抵抗〔K/W〕（波-4）、音源からの距離（波-7）、曲率半径（波-12・13）。比気体定数は $R_{\mathrm s}$（波-9） & 電気抵抗（電-4）、気体定数（熱-3） & 単位：Ω、K/W、m \\
$\delta$ & 対数減衰率（波-1）、プリズムの偏角（波-12）。静たわみは $x_{\mathrm{st}}$ & デルタ関数 $\delta(t)$（9-6）、道で変わる小さな量（熱-2） & 単位なし、角 \\
$\lambda$ & 波長〔m〕（第2〜4部）、固有値 $\omega^2$（波-3 だけ） & 管摩擦係数（力-14）、固有値（7-8） & 単位：m、1/s$^2$ \\
$\theta$ & 角（ふりこ・入射角など）、セルシウス温度〔℃〕（波-4・9） & セルシウス温度（熱-1）、角（4-2） & 単位：rad・°、℃ \\
$p$、$P$ & 音圧〔Pa〕（波-7・9）、確率 $p_i$（波-11）、仕事率・電力 $P$〔W〕（波-7・11） & 圧力（熱）、運動量 $\vec p$（力-3） & 単位：Pa、W \\
$r$ & 反射係数（波-7・10・13）、距離（波-9）、屈折角（波-13）、目標の値（波-4）、不つり合いの半径 $r_{\mathrm u}$（波-2） & 位置（力-1）、相関係数（3-14） & 反射係数は単位なし \\
$a$、$d$ & スリットの幅 $a$・間かく $d$（波-8・13）、物までの距離 $a$（波-12）、加速度 $a$（波-10）、圧電定数 $d_{33}$（波-10）、外乱 $d$（波-4） & 加速度（力-1）、音速 $a$（熱-4） & 単位：m、m/s$^2$、C/N \\
$v$ & 速度、物質の中の光の速さ $v=c/n$（第4部）、音源・観測者の速さ $v_{\mathrm s}$・$v_{\mathrm o}$（波-8） & 速度（力-1）、比体積（熱-8） & 単位：m/s \\
$\alpha$ & 吸音率（波-9）、吸収の定数〔1/m〕（波-7）、円すいの半角（波-8）、プリズムの頂角（波-12） & 線膨張係数（力-8・熱-1） & 吸音率は単位なし \\
\end{longtblr}\par}

\phbkhead{この巻の約束（記号と向き）}
{\small
\begin{itemize}[leftmargin=1.2\zw, itemsep=2pt, topsep=1pt, label={\color{bkVol}●}]
\item \textbf{$c$ の3つの意味}：第1部は減衰係数、第2・3部は波の速さ、第4部は光速です。同じ章では2つの意味にしません。高校の教科書は、波の速さを $v$、音の速さを $V$ と書きます。
\item \textbf{振幅}：振幅は $A$、複素振幅は $X$（大きさ $\lvert X\rvert$ が振幅、偏角がおくれ。数学 6-13 と同じ）。おくれは負の角で書きます。虚数単位は $j$ です（数学 6-8）。
\item \textbf{振動数}：振動数 $f$〔Hz〕と角振動数 $\omega$〔rad/s〕を書き分けます（$\omega=2\pi f$）。回る機械は「毎分1800回転」のように書き、$\qty{30}{Hz}$ と同じです。
\item \textbf{反射係数の向き}：波-7 の $r$ は変位（ずれ）や速さで数え、音響の本は音圧で数えるので、符号が逆になります（波-10 は音圧で数えた $(Z_2-Z_1)/(Z_2+Z_1)$）。はね返るエネルギーの割合 $r^2$ は、どちらでも同じです。
\item \textbf{ドップラー効果の向き}：音源・観測者の速さ $v_{\mathrm s}$・$v_{\mathrm o}$ は、音源から観測者へ向かう向きを正にします（波-8）。
\item \textbf{デシベル}：いつも基準を書きます。音圧は\qty{20}{\micro Pa}、音の強さは $10^{-12}$\,W/m$^2$、音響パワーは $10^{-12}$\,W です（波-9）。強さ・パワーの比は $10\log_{10}$、音圧・振幅の比は $20\log_{10}$ です。
\end{itemize}\par}

%% ---------------- 音速・屈折率・物性値 ----------------
\bkbacktitle{音速・屈折率・物性値}
{\small この巻で使った値です。本文の例題の値は、この表と同じ丸めです。物性値は温度や材料の種類で変わるので、設計では便覧・規格の値を使います。出典の欄の「波-7」などは、その値を使った章です。\phyk は、規格・公式資料の本文でも、それを引く2つ以上の資料でも、まだ確かめられていない値です。\par}

\phbkhead{1　音の速さ・密度・音響インピーダンス}
{\footnotesize
\begin{longtblr}{colspec={Q[l,m,wd=40mm] Q[c,m] Q[c,m] Q[c,m] Q[l,m,wd=24mm]}, width=\linewidth, rowhead=1, row{1}={bg=bkVol!14, font=\sffamily\bfseries\footnotesize}, hlines={0.4pt,bkLine}, rowsep=1.4pt, colsep=3.5pt, colsep=2.5pt}
物 & {密度 $\rho$\\〔kg/m$^3$\,〕} & {音速 $c$\\〔m/s〕} & {$Z=\rho c$\\〔Pa·s/m〕} & 出典 \\
空気（\qty{0}{\celsius}、1気圧） & 1.293 & 331 & 429 & ア（波-9） \\
空気（\qty{20}{\celsius}、1気圧） & 1.204 & 343 & 413 & ア（熱-4・波-7・波-9） \\
ヘリウム（\qty{20}{\celsius}、1気圧） & 0.1663 & 1008 & 168 & イ（熱-4・波-9） \\
水（\qty{8}{\celsius}） & 999.9 & 1439 & $1.44\times10^6$ & イ（波-9） \\
水（\qty{20}{\celsius}） & 998.2 & 1482（約1480） & $1.48\times10^6$ & イ（波-7・波-9） \\
海水 & ― & 約1500 & ― & ウ（波-9・波-10） \\
体の軟らかい組織 & 約1060 & 約1540 & $1.63\times10^6$ & エ（波-7・8・10） \\
超音波検査のゼリー & ― & 約1530 & 約 $1.5\times10^6$ & オ（波-7） \\
骨 & ― & ― & 約 $7.8\times10^6$ & エ（波-7） \\
圧電セラミックス & ― & ― & 約 $3\times10^7$ & カ（波-7） \\
鋼（細い棒の縦波） & 7850 & 約5100 & ― & キ（波-9・力-16） \\
鋼（広い固体の縦波） & 7850 & 約5900 & 約 $4.6\times10^7$ & キ（波-7・9・10） \\
鋼（横波） & 7850 & 約3200 & ― & キ（波-9） \\
\end{longtblr}\par}
{\footnotesize\color{bkGray}出典　ア：CoolProp 8.0（空気の状態式）で計算（343.3\,m/s・1.2046\,kg/m$^3$）。理想気体の式 $\sqrt{\gamma R_{\mathrm s}T}$（$\gamma=1.4$、$R_{\mathrm s}=\qty{287}{J/(kg.K)}$）で 331.3・343.2（熱-4 ⑥）。イ：NIST Chemistry WebBook（流体の熱物性。水 \qty{20}{\celsius}で 998.21\,kg/m$^3$・1482.35\,m/s、ヘリウム 1007.9\,m/s）。ウ：海の音の解説（DOSITS「How fast does sound travel?」ほか）。エ：医用超音波の講義資料（スタンフォード大学 RAD225 ほか）。オ：アメリカ食品医薬品局の届け出（510(k) K102637。1530\,m/s・$1.505\times10^6$\,Pa·s/m）。カ：特許文献の値（$3.0$〜$3.4\times10^7$）。キ：$E=\qty{205}{GPa}$・$\nu=0.3$・$G=\qty{79}{GPa}$ から計算（棒 $\sqrt{E/\rho}=5110$、縦 $\sqrt{(K+\frac43G)/\rho}=5930$、横 $\sqrt{G/\rho}=3170$）。超音波探傷の音速表の 5890\,m/s（炭素鋼）と合う。\par
空気の音速は、\qty{0}{\celsius}から\qty{1}{K}上がるごとに約\qty{0.6}{m/s}速くなる（$331.5+0.6\theta$、波-9 ⑤）。固い物ほど速く、軽い物ほど速い：$c=\sqrt{\text{縮みにくさ}/\rho}$（波-5 ⑤）。\par}

\phbkhead{2　弾性の定数（室温）}
{\footnotesize
\begin{longtblr}{colspec={Q[l,m,wd=25mm] Q[c,m,wd=17mm] Q[c,m,wd=14mm] Q[c,m,wd=12mm] Q[c,m,wd=14mm] Q[c,m,wd=14mm] X[l,m]}, width=\linewidth, rowhead=1, row{1}={bg=bkVol!14, font=\sffamily\bfseries\footnotesize}, hlines={0.4pt,bkLine}, rowsep=1.4pt, colsep=3.5pt, colsep=2pt}
物 & {密度 $\rho$\\〔kg/m$^3$\,〕} & {ヤング率\\$E$\\〔GPa〕} & {横弾性\\係数 $G$\\〔GPa〕} & {ポアソン\\比 $\nu$} & {体積弾性\\率 $K$\\〔GPa〕} & 出典 \\
鋼 & 7850 & 205 & 79 & 0.3 & 約171 & 第1巻の巻末と同じ（力-8・力-11）。$K$ は計算（波-9） \\
板ガラス（ソーダ石灰） & 約2500 & 72 & ― & 0.23 & ― & ガラスの会社の技術資料（Pilkington ATS-129。波-9） \\
水（\qty{20}{\celsius}） & 998 & ― & ― & ― & 約2.2 & 力-8・力-12（$\rho c^2$ で 2.19） \\
空気（\qty{20}{\celsius}、断熱） & 1.204 & ― & ― & ― & 0.000\,14 & $\gamma p$（熱-3 ⑧） \\
\end{longtblr}\par}
{\footnotesize\color{bkGray}体積弾性率は「押し縮めにくさ」、ヤング率は「引きのばしにくさ」、横弾性係数は「ずらしにくさ」。$G=E/(2(1+\nu))$、$K=E/(3(1-2\nu))$（力-8・力-11）。空気の $\gamma p=1.4\times\qty{101325}{Pa}\approx\qty{0.14}{MPa}$。\par}

\phbkhead{3　屈折率（黄色の光 約\qty{589}{nm}、\qty{20}{\celsius}）}
{\footnotesize
\begin{longtblr}{colspec={Q[l,m,wd=38mm] Q[c,m,wd=28mm] Q[c,m,wd=16mm] X[l,m]}, width=\linewidth, rowhead=1, row{1}={bg=bkVol!14, font=\sffamily\bfseries\footnotesize}, hlines={0.4pt,bkLine}, rowsep=1.4pt, colsep=3.5pt}
物 & 屈折率 $n$ & 臨界角 & 出典・注意 \\
空気（1気圧） & 1.0003 & ― & Ciddor（1996）の式で\qty{15}{\celsius} 1.000\,277。\qty{20}{\celsius}・\qty{633}{nm}で 1.000\,27（波-12・波-13） \\
水 & 1.333（本文 1.33） & \ang{49} & Tilton・Taylor の測定値を引く表。Daimon・Masumura（2007）の式で 1.3333（波-12・波-13） \\
水（\qty{405}{nm}・\qty{656}{nm}） & 1.343・1.331 & ― & 同じ表。色で曲がり方がちがう（分散、波-12 ⑧）。Daimon・Masumura の式で 1.3431・1.3315 \\
ガラス（クラウンガラス BK7） & 1.5168（本文 約1.5・1.52） & \ang{42} & SCHOTT の N-BK7 のデータシート（\qty{587.6}{nm}）。臨界角は $n=1.5$ で（波-12・波-13） \\
ダイヤモンド & 2.417（本文 2.42） & \ang{24} & 宝石の屈折率の表。Peter（1923）の式で 2.417（波-12） \\
フッ化マグネシウム（反射防止の膜） & 1.38 & ― & 光学部品の会社の資料（\qty{550}{nm}の膜）。Dodge（1984）の式で結晶は 1.378（波-7・波-13） \\
光ファイバ（コア・クラッド） & 1.48・1.46 & \ang{80.6} & 講義資料の計算例。開口数 0.24、受け入れ角 約\ang{14}（波-12） \\
\end{longtblr}\par}
{\footnotesize\color{bkGray}臨界角 $\theta_{\mathrm c}$ は $\sin\theta_{\mathrm c}=1/n$（空気との境。光ファイバはクラッドとの境で $1.46/1.48$）。ガラスに垂直に当たる光の反射は約4\,\%（$((n-1)/(n+1))^2$、波-7 ⑨）。\par}

\phbkhead{4　光と音の波長・振動数の目安}
{\footnotesize
\begin{longtblr}{colspec={Q[l,m,wd=40mm] Q[l,m,wd=42mm] X[l,m]}, width=\linewidth, rowhead=1, row{1}={bg=bkVol!14, font=\sffamily\bfseries\footnotesize}, hlines={0.4pt,bkLine}, rowsep=1.4pt, colsep=3.5pt}
名前 & 値 & 出典・注意 \\
可視光（目に見える光） & 約 \qtyrange{380}{780}{nm} & 国際照明委員会 CIE S 017:2020（下の限界 360〜400\,nm、上の限界 760〜830\,nm で、決まった境はない）。熱-13・波-12・波-13 \\
ナトリウムの黄色の光（D 線） & \qty{589.0}{nm}・\qty{589.6}{nm} & NIST の原子スペクトルの表。屈折率の表の「589\,nm」 \\
赤いレーザー（ヘリウムネオン） & \qty{633}{nm}（\qty{632.8}{nm}） & NIST の原子スペクトルの表（ネオン）。波-13 \\
CD・DVD・BD の読み取りの光 & \qty{780}{nm}・\qty{650}{nm}・\qty{405}{nm} & 光ディスクの比較表・特許の文書（波-13） \\
CD・DVD・BD のトラックの間かく & \qty{1.6}{\micro m}・\qty{0.74}{\micro m}・\qty{0.32}{\micro m} & 同じ（波-13） \\
人に聞こえる音 & 約\qty{20}{Hz}〜\qty{20}{kHz}（若い人） & 用語解説・音響の会社の資料。\qty{20}{kHz}より上が超音波（波-9・波-11） \\
救急車のサイレン & \qty{960}{Hz}・\qty{770}{Hz}（各 0.65 秒） & 東京消防庁の資料（波-8） \\
コウモリの超音波 & 国内の小型のコウモリ 約\qtyrange{10}{120}{kHz}、アブラコウモリ 約\qty{40}{kHz} & 国土技術政策総合研究所の資料・日本心理学会の資料（波-8・波-10） \\
超音波の機械 & 駐車の支援 \qtyrange{40}{70}{kHz}、魚群探知機 \qty{50}{kHz}・\qty{200}{kHz} & 部品の会社・機械の取扱説明書（波-10） \\
医用の超音波・鋼の溶接部の探傷 & 約\qtyrange{2}{10}{MHz}・約\qtyrange{2}{5}{MHz} & 医用超音波の講義資料、探傷の会社の資料（波-10） \\
速度をはかるレーダー & \qty{24}{GHz}帯（本文は\qty{24.15}{GHz}とする） & 総務省の無線局の審査基準（波-8） \\
水晶時計の振動子 & \qty{32768}{Hz}（$2^{15}$） & 計算（波-10） \\
\end{longtblr}\par}

\phbkhead{5　減衰比・Q 値・ゆれの振動数の目安}
{\footnotesize
\begin{longtblr}{colspec={Q[l,m,wd=40mm] Q[l,m,wd=42mm] X[l,m]}, width=\linewidth, rowhead=1, row{1}={bg=bkVol!14, font=\sffamily\bfseries\footnotesize}, hlines={0.4pt,bkLine}, rowsep=1.4pt, colsep=3.5pt}
名前 & 値 & 出典・注意 \\
鋼の部品1つ（つなぎ目がない） & 減衰比 0.001〜0.01 & ボルト・溶接の減衰の技術資料、米国原子力規制委員会の指針で照らし合わせた（波-1） \\
高い建物（鉄骨） & 減衰比 0.01〜0.02 & 強震観測の論文・規制委員会の資料（波-1） \\
建物（鉄筋コンクリート） & 減衰比 0.02〜0.05 & 同じ（波-1） \\
設計で使う建物の減衰比 & 鉄骨造 約0.02、鉄筋コンクリート造 約0.03 & 日本建築センターの高層の評定の集計を引く資料（波-3）。実測では 0.015〜0.03 \\
車のサスペンション & 減衰比 0.2〜0.4 & Gillespie『Fundamentals of Vehicle Dynamics』5章を引く講義資料（波-1・波-2） \\
制御でよく使う2次おくれ & 減衰比 0.5〜0.7（行き過ぎ 16〜5\,\%） & 計算（波-4） \\
Q 値 & 建物 約10〜50、音さ 数千、時計の水晶振動子 1万以上 & $Q\approx1/(2\zeta)$。音さの論文（Am.\ J.\ Phys.）、NIST の水晶の資料（波-2） \\
乗用車の上下のゆれ & 約\qtyrange{1}{1.5}{Hz} & 車両運動の論文・解説（波-1） \\
人が歩く足の運び & 約\qty{2}{Hz}（横にゆする力は約\qtyrange{0.7}{1.2}{Hz}） & 歩道橋の設計の総説（波-1・波-2） \\
ミレニアム橋（ロンドン）の横ゆれ & \qty{0.48}{Hz}・\qty{0.95}{Hz}（中央）、\qty{0.77}{Hz}（南） & Ingenia の記事・構造の論文（波-2） \\
高い建物の1次の周期 & 階数×0.1 秒くらい。設計の略算は 高さ〔m〕×(0.02〜0.03) 秒 & 気象庁「長周期地震動について」ほか（波-1）。略算 $T=h(0.02+0.01\alpha)$ は建設省告示第1793号（波-3） \\
動吸振器 & 質量比 1〜5\,\%。台北101 は直径\qty{5.5}{m}・\qty{660}{t}の鋼の球 & 制振装置の会社の資料、ギネス世界記録ほか（波-3） \\
円柱の後ろの渦（ストローハル数） & 約0.2 & 力-15（波-2） \\
回る機械 & 4極モーター 毎分約1500回転（\qty{50}{Hz}）・1800回転（\qty{60}{Hz}） & 同期速度 $120f/\text{極の数}$ から（波-2・波-10） \\
家庭の洗濯機の脱水 & 毎分400〜1100回転ほど & 特許文献・海外の販売店の解説。国内のメーカーの仕様では確かめていない\phyk（波-2） \\
\end{longtblr}\par}

\phbkhead{6　A 特性の補正（騒音計）}
{\footnotesize
\begin{longtblr}{colspec={Q[l,m,wd=24mm] *{11}{Q[c,m]}}, width=\linewidth, rowhead=1, row{1}={bg=bkVol!14, font=\sffamily\bfseries\footnotesize}, hlines={0.4pt,bkLine}, rowsep=1.4pt, colsep=3.5pt, colsep=1.6pt}
振動数〔Hz〕 & 31.5 & 63 & 100 & 125 & 250 & 500 & 1000 & 2000 & 4000 & 8000 & \num{16000} \\
補正〔dB〕 & $-39.4$ & $-26.2$ & $-19.1$ & $-16.1$ & $-8.6$ & $-3.2$ & 0 & $+1.2$ & $+1.0$ & $-1.1$ & $-6.6$ \\
\end{longtblr}\par}
{\footnotesize\color{bkGray}IEC 61672-1:2013（JIS C 1509-1:2017）の表の値。規格の式（20.6・107.7・737.9・\qty{12194}{Hz}、波-9 ⑧）で計算した値と小数第1位まで同じ。はかった音圧レベルに、この補正をたすと dB(A) になる。\par}

\phbkhead{7　音の大きさと騒音の決まり}
{\footnotesize
\begin{longtblr}{colspec={Q[l,m,wd=40mm] Q[l,m,wd=42mm] X[l,m]}, width=\linewidth, rowhead=1, row{1}={bg=bkVol!14, font=\sffamily\bfseries\footnotesize}, hlines={0.4pt,bkLine}, rowsep=1.4pt, colsep=3.5pt}
名前 & 値 & 出典・注意 \\
基準の値（0\,dB） & 音圧 \qty{20}{\micro Pa}、音の強さ $10^{-12}$\,W/m$^2$、音響パワー $10^{-12}$\,W & 音響の講義資料（MIT・大学の資料）（波-9） \\
聞こえ始める音・会話・ガード下 & \qty{20}{\micro Pa}（0\,dB）・\qty{0.02}{Pa}（60\,dB）・\qty{2}{Pa}（100\,dB） & 波-9 \\
耳が痛くなるほどの音 & 約\qtyrange{20}{60}{Pa}（120〜130\,dB） & 同じ \\
ささやき・図書館・工場 & 約30\,dB・40\,dB・90\,dB & 数学 3-12⑦と同じ数（波-9） \\
騒音に係る環境基準（等価騒音レベル） & 住宅の地域（A・B）昼55\,dB・夜45\,dB 以下、特に静かな地域（AA）50・40、商業・工業もある地域（C）60・50、幹線道路に近い所 70・65 & 環境庁告示第64号（環境省の本文で確かめた）。A 特性ではかる（波-9） \\
働く人の耳を守る目安 & 等価騒音レベル 85\,dB 以上の作業場で対策 & 厚生労働省「騒音障害防止のためのガイドライン」（令和5年改訂）（波-9） \\
救急車のサイレンの大きさ & 前へ\qty{20}{m}の所で 90\,dB 以上 120\,dB 以下 & 道路運送車両の保安基準の細目を定める告示（国土交通省の本文で確かめた）（波-8） \\
板ガラス\qty{5}{mm}の透過損失（\qty{500}{Hz}） & 約30\,dB & 質量則で計算（波-9） \\
残響時間 & 教室 0.6 秒以下、コンサートホール 約2秒前後 & ANSI S12.60（教室）、音響の論文（ホール）（波-9） \\
コンクリートの吸音率（\qty{500}{Hz}） & 0.02〜0.04 & 吸音材の会社の表（波-9） \\
\end{longtblr}\par}

\phbkhead{8　実機の値と目安（計測・信号・制御・光）}
{\footnotesize
\begin{longtblr}{colspec={Q[l,m,wd=40mm] Q[l,m,wd=42mm] X[l,m]}, width=\linewidth, rowhead=1, row{1}={bg=bkVol!14, font=\sffamily\bfseries\footnotesize}, hlines={0.4pt,bkLine}, rowsep=1.4pt, colsep=3.5pt}
名前 & 値 & 出典・注意 \\
体の中の超音波の吸収 & 約 \mbox{0.5\,dB/(cm·MHz)} & 医用超音波の講義資料（波-7・波-10） \\
圧電定数 $d_{33}$ & PZT 約\qtyrange{300}{700}{pC/N}、水晶 約\qty{2.3}{pC/N} & IEEE の解説・論文（波-10） \\
圧電式の加速度計 & 固有振動数の約3分の1まで使う。上限 約\qtyrange{5}{50}{kHz} & 計測器の会社の解説（波-10） \\
CD の音 & 1秒に\num{44100}回・16ビット・2チャンネル（\qty{1411.2}{kbit/s}） & CD の規格（IEC 60908）。数学 9-4・9-12 と同じ（波-11） \\
固定電話のディジタルの声 & 1秒に8000回・8ビット（\qty{64}{kbit/s}） & ITU-T G.711（波-11） \\
電話の線のモデム & 最大 \qty{33.6}{kbit/s} & ITU-T V.34（波-11） \\
量子化の S/N 比 & 8ビット 約50\,dB、16ビット 約98\,dB & $6.02b+1.76$\,dB から（波-11） \\
制御の余裕の目安 & ゲイン余裕 6〜12\,dB、位相余裕 \ang{30}〜\ang{60} & 制御の講義資料・半導体の会社の資料（波-4） \\
人の反応のおくれ（目で見て手を動かす） & 約0.2 秒 & 眼科の論文（平均 208〜224\,ms）（波-4） \\
ドローンの姿勢の制御 & 1秒に数百回以上（PID の計算は 1〜\qty{8}{kHz}） & 制御ソフトの説明書（波-4） \\
事務室の照度 & \qty{750}{lx} & JIS Z 9110:2010 の値を引く電力会社・照明の会社の資料（波-12） \\
晴れた日の直射日光の照度 & 約10万\,lx & 照明の会社の資料・照明学会誌（波-12） \\
LED 電球（60形相当）の光束 & \qty{810}{lm} & 照明の会社の発表（日本電球工業会のガイドライン）（波-12） \\
空気だけによる散乱（\qty{550}{nm}） & 光学的な厚さ 0.097。地平線の向きは空気の量が約38倍 & Bodhaine ほか（1999）、Kasten・Young（1989）の式（波-13） \\
重力波の検出器 & LIGO の腕 \qty{4}{km}（のび縮みの割合 $1.0\times10^{-21}$）、KAGRA の腕 \qty{3}{km} & LIGO の論文（Phys.\ Rev.\ Lett.\ 116, 061102）、KAGRA の論文（波-13） \\
\end{longtblr}\par}

\phbkhead{9　この巻で使う定数}
{\footnotesize
\begin{longtblr}{colspec={Q[l,m,wd=40mm] Q[l,m,wd=46mm] X[l,m]}, width=\linewidth, rowhead=1, row{1}={bg=bkVol!14, font=\sffamily\bfseries\footnotesize}, hlines={0.4pt,bkLine}, rowsep=1.4pt, colsep=3.5pt}
名前 & 値 & 出典・注意 \\
真空の光速 $c$ & \qty{299792458}{m/s}（決めた値） & SI。NIST の CODATA の値で確かめた。本文では約 $3.00\times10^8$\,m/s（波-12・波-13） \\
1\,m の決め方 & 光が真空を $1/\num{299792458}$ 秒に進む長さ & 1983年に決めた（BIPM・NIST）（波-12・波-13） \\
標準重力加速度 & \qty{9.80665}{m/s^2}（決めた値） & NIST の CODATA の値。本文では 9.8（波-1） \\
標準大気圧 & \qty{101325}{Pa}（決めた値） & 同じ（波-9） \\
\qty{0}{\celsius} & \qty{273.15}{K}（定義） & $\theta$/℃＝$T$/K$-273.15$（熱-1・波-9） \\
空気の比熱比・比気体定数 & $\gamma=1.4$、$R_{\mathrm s}=\qty{287}{J/(kg.K)}$ & 熱-3 と同じ（波-9） \\
ボルツマン定数 $k_{\mathrm B}$ & \qty{1.380649e-23}{J/K}（決めた値） & SI（熱-15・波-11） \\
太陽が出す光の仕事率 & \qty{3.828e26}{W} & IAU 2015 Resolution B3（波-7） \\
太陽定数（地球の所の強さ） & \qty{1361}{W/m^2} & 同じ（波-7。熱-13 と同じ） \\
\end{longtblr}\par}

\phbkhead{10　単位と換算}
{\footnotesize
\begin{longtblr}{colspec={Q[l,m,wd=34mm] Q[l,m,wd=34mm] X[l,m]}, width=\linewidth, rowhead=1, row{1}={bg=bkVol!14, font=\sffamily\bfseries\footnotesize}, hlines={0.4pt,bkLine}, rowsep=1.4pt, colsep=3.5pt}
量 & SI の単位 & ほかの単位との関係（単位つきの1で換算する。2-10） \\
振動数 & Hz（ヘルツ）＝1/s & 毎分 $N$ 回転＝$N/60$\,Hz（毎分1800回転＝\qty{30}{Hz}） \\
角振動数 & rad/s & $\omega=2\pi f$。\qty{1}{Hz}は約\qty{6.28}{rad/s} \\
波長 & m & 1\,nm＝$10^{-9}$\,m、1\,μm＝$10^{-6}$\,m \\
レベル & dB（デシベル） & 強さ・パワー：$10\log_{10}$（比）。音圧・電圧・振幅：$20\log_{10}$（比）。2倍で約3\,dB・約6\,dB \\
音響インピーダンス & Pa·s/m & 専門書の rayl（レイル）は Pa·s/m と同じ。1\,MRayl＝$10^6$\,Pa·s/m \\
光の量 & lm・cd・lx & 照度 1\,lx＝1\,lm/m$^2$。光度 1\,cd は、ある向きの立体角1\,sr あたり 1\,lm \\
情報 & ビット・ビット/s & 1\,kbit/s＝1000 ビット/s。8ビット＝1バイト \\
角 & rad & 1°＝$\pi/180$\,rad。1分（角）＝1°の60分の1≈$2.9\times10^{-4}$\,rad \\
\end{longtblr}\par}
{\footnotesize\color{bkGray}接頭語：G（ギガ、$10^9$）、M（メガ、$10^6$）、k（キロ、$10^3$）、m（ミリ、$10^{-3}$）、μ（マイクロ、$10^{-6}$）、n（ナノ、$10^{-9}$）、p（ピコ、$10^{-12}$）。\par}

%% ---------------- なぜ？の道すじ ----------------
\bkbacktitle{なぜ？の道すじ}
{\small 気になる「なぜ？」から、読む章の順番をたどる路線図です。札の横の数はこの巻のページです。ほかの巻の札（力・熱・電・原）は、物理シリーズの『力学』『熱学』『電磁気学』『原子』の章です。使う数学の章は、下の小さな札です。\par}
\vspace{1mm}
\phwhyroute{地震のゆれ、楽器の音、光。どこが同じ？}{波-0}{波-0,波-1,波-5,波-12}{6-6}
\phwhyroute{ブランコの1往復の時間は、何で決まる？}{波-1}{力-2,波-1}{5-17,6-6}
\phwhyroute{洗濯機は、なぜ回り始めにガタガタゆれることがある？}{波-2}{波-1,波-2}{6-7,6-13}
\phwhyroute{高いビルの屋上の重りは、どうやってゆれをおさえる？}{波-3}{波-1,波-2,波-3}{6-14,7-8}
\phwhyroute{エアコンは、どうやって部屋をちょうどよい温度に保つ？}{波-4}{熱-11,熱-12,波-4}{6-4,9-8}
\phwhyroute{スタジアムのウェーブでは、何が進んでいる？}{波-5}{波-1,波-5}{8-13}
\phwhyroute{ギターの弦は、なぜおさえる場所で音の高さが変わる？}{波-6}{波-5,波-6}{9-1}
\phwhyroute{超音波の検査で、なぜおなかにゼリーをぬる？}{波-7}{波-5,波-7,波-10}{3-12}
\phwhyroute{救急車のサイレンは、なぜ通りすぎると低く聞こえる？}{波-8}{波-5,波-8}{5-17}
\phwhyroute{雷が光ってから、なぜ少しおくれて音が聞こえる？}{波-9}{熱-4,波-9}{}
\phwhyroute{コウモリは、暗やみでどうやってものの位置を知る？}{波-10}{波-8,波-9,波-10}{}
\phwhyroute{動画の車輪が、逆に回って見えることがあるのはなぜ？}{波-11}{波-11}{9-4}
\phwhyroute{水に入れたストローは、なぜ折れて見える？}{波-12}{波-8,波-12}{4-6,5-7}
\phwhyroute{空は、なぜ青い？　夕焼けは、なぜ赤い？}{波-13}{波-1,波-2,波-5,波-13,電-15}{6-13}
\phwhyroute{スマホの声は、どうやって電波にのせて送る？}{波-14}{波-11,波-14,電-15}{9-12}
\phwhyroute{ふりこ時計は、なぜ正確に時をきざめる？}{波-1}{力-2,波-1}{5-17}
\phwhyroute{歩道橋が、人が歩くと横にゆれることがあるのはなぜ？}{波-2}{波-1,波-2}{6-13}
\phwhyroute{ドローンは、どうやって空中で同じ姿勢を保つ？}{波-4}{波-1,波-4}{9-8,9-9}
\phwhyroute{ヘリウムを吸って話すと、なぜ声が変わる？}{波-9}{熱-4,波-6,波-9}{}
\phwhyroute{線路に耳をあてると、なぜ遠くの汽車の音が早く聞こえる？}{波-9}{力-8,波-5,波-9}{}
\phwhyroute{同じ音を2台にしても、なぜ2倍の大きさに聞こえない？}{波-9}{波-7,波-9}{3-12}
\phwhyroute{重い壁ほど、なぜ外の音が聞こえにくい？}{波-9}{波-2,波-7,波-9}{3-12}
\phwhyroute{地震のとき、なぜ先に小さくゆれて、あとから大きくゆれる？}{波-9}{波-5,波-9}{}
\phwhyroute{広い部屋で手をたたくと、なぜ音がしばらく残る？}{波-9}{波-7,波-9}{3-12}
\phwhyroute{海岸の波は、なぜ岸に平行になる？}{波-8}{波-5,波-8}{}
\phwhyroute{音より速い飛行機が通ると、なぜ「ドーン」と聞こえる？}{波-8}{熱-4,波-8}{4-6}
\phwhyroute{雑音を消すヘッドホンは、どうやって音を消す？}{波-8}{波-8,波-11}{4-9}
\phwhyroute{鉄の中の見えない傷を、どうやって見つける？}{波-10}{波-7,波-10}{}
\phwhyroute{機械のゆれから、どうやって故障を見つける？}{波-10}{波-2,波-10,波-11}{9-4}
\phwhyroute{魚群探知機は、どうやって海の中の魚を見つける？}{波-10}{波-9,波-10}{}
\phwhyroute{CD は、なぜ1秒に\num{44100}回も音をはかる？}{波-11}{波-11}{9-4}
\phwhyroute{光ファイバは、なぜ曲がっていても光を運べる？}{波-12}{波-12}{4-6}
\phwhyroute{虹は、なぜ色が分かれて見える？}{波-12}{波-12,波-13}{}
\phwhyroute{シャボン玉は、なぜ色がついて見える？}{波-13}{波-7,波-13}{4-9}
\phwhyroute{大きな望遠鏡ほど、なぜ近い2つの星を見分けられる？}{波-13}{波-8,波-12,波-13}{}
\phwhyroute{偏光サングラスをかけると、なぜ水面のぎらぎらが消える？}{波-13}{波-12,波-13}{4-6}
\phwhyroute{空はなぜ紫ではない？　雲はなぜ白い？}{波-13}{波-2,波-13}{}

%% ---------------- チェック表 ----------------
\bkbacktitle{できるようになったことチェック表}
{\small 各章の「ふりかえり」を集めた表です。章を学び終えたら1回め、次の巻に進む前に2回めのチェックをしよう。自信がない所は、ページの数字の所へもどろう。\par}\vspace{2mm}
% 自動生成（tools/mkchecklist.py p3）。各章の「ふりかえり」の項目を集めたもの（第13回）
\begingroup\footnotesize
\def\bkcb{\tikz[baseline=-0.5ex]\draw[\bkTF, line width=0.7pt, rounded corners=0.5pt] (0,-0.3em) rectangle (0.8em,0.5em);}
\begin{longtblr}{colspec={Q[c,m,wd=11mm] X[l,m] Q[c,m,wd=10mm] Q[c,m,wd=9mm] Q[c,m,wd=9mm]}, width=\linewidth, rowhead=1, colsep=3pt,
  row{1}={bg=bkVol!14, font=\sffamily\bfseries\scriptsize}, hlines={0.4pt,bkLine}, rowsep=0.7pt}
章 & できるようになったこと & ページ & 1回め & 2回め \\
{\sffamily\bfseries\color{\bkTT}波-0} & ゆれの周期と振動数、波の波長と速さを、身近な例で説明できる。 & \pageref{ph:波-0:sec2} & \bkcb & \bkcb \\
 & 何を無視したかを書いてから、波が届くまでの時間や距離を、7つの段で求められる。 & \pageref{ph:波-0:ex3} & \bkcb & \bkcb \\
 & 身近な「なぜ？」の答えが、この巻のどの章にあるかを探せる。 & \pageref{ph:波-0:sec6} & \bkcb & \bkcb \\
{\sffamily\bfseries\color{\bkTT}波-1} & ふりこの1往復の時間が、ふりこの長さで決まり、おもりの重さにはよらないことを、実験で確かめられる。 & \pageref{ph:波-1:sec3} & \bkcb & \bkcb \\
 & ばねとおもり・ふりこの運動方程式から、単振動の式と周期を導き、エネルギーで速さを求められる。 & \pageref{ph:波-1:sec2} & \bkcb & \bkcb \\
 & 等価ばね定数・静たわみ・実体ふりこ・ねじり振動で、いろいろな物の固有振動数を見積もれる。 & \pageref{ph:波-1:sec4} & \bkcb & \bkcb \\
 & ゆれの記録から減衰比を読み、力学・電気・熱・流体の相似の表で、同じ形の式を見分けられる。 & \pageref{ph:波-1:sec7} & \bkcb & \bkcb \\
{\sffamily\bfseries\color{\bkTT}波-2} & 共振（共鳴）を、ゆする速さとゆれの大きさの関係で説明できる。 & \pageref{ph:波-2:sec1} & \bkcb & \bkcb \\
 & 表計算でばねをゆすって、共振曲線とゆれのおくれを読みとれる。 & \pageref{ph:波-2:sec3} & \bkcb & \bkcb \\
 & 回る物の不つり合い・でこぼこ道・風の渦が、どれだけの力で、1秒に何回ゆするかを求められる。 & \pageref{ph:波-2:kazoeru} & \bkcb & \bkcb \\
 & 共振の倍率・Q 値・伝達率を求め、防振ばねを選べる。 & \pageref{ph:波-2:sec5} & \bkcb & \bkcb \\
 & 不つり合い・回る軸・道のでこぼこ・風の渦によるゆれを見積もり、ボード線図で読める。 & \pageref{ph:波-2:sec7} & \bkcb & \bkcb \\
{\sffamily\bfseries\color{\bkTT}波-3} & ゆれが行き来するわけを、2つのモードのたし算で説明できる。 & \pageref{ph:波-3:sec2} & \bkcb & \bkcb \\
 & 表計算で2自由度のゆれと、動吸振器の効きめを確かめられる。 & \pageref{ph:波-3:sec3} & \bkcb & \bkcb \\
 & ラグランジュの式で、エネルギーから運動方程式を立てられる。 & \pageref{ph:波-3:sec4} & \bkcb & \bkcb \\
 & 質量行列のある固有値問題で、モードを求められる。 & \pageref{ph:波-3:sec5} & \bkcb & \bkcb \\
 & 動吸振器の同調と、モードの直交・モード解析を説明できる。 & \pageref{ph:波-3:sec6} & \bkcb & \bkcb \\
{\sffamily\bfseries\color{\bkTT}波-4} & フィードバックのしくみを、ブロック線図で説明できる。 & \pageref{ph:波-4:sec2} & \bkcb & \bkcb \\
 & 表計算で部屋の温度の比例・積分の制御を追える。 & \pageref{ph:波-4:sec3} & \bkcb & \bkcb \\
 & 物理の法則から伝達関数を作り、応答の速さを読める。 & \pageref{ph:波-4:sec4} & \bkcb & \bkcb \\
 & 残るずれと PID の働きを、最終値の定理で説明できる。 & \pageref{ph:波-4:sec6} & \bkcb & \bkcb \\
 & ラウスの表で、安定かどうかを判定できる。 & \pageref{ph:波-4:sec8} & \bkcb & \bkcb \\
 & ゲイン余裕・位相余裕と状態フィードバックで、安定の度合いを考えられる。 & \pageref{ph:波-4:sec9} & \bkcb & \bkcb \\
{\sffamily\bfseries\color{\bkTT}波-5} & 波で進むのは「ゆれの形」だと、横波・縦波で説明できる。 & \pageref{ph:波-5:sec1} & \bkcb & \bkcb \\
 & 振幅・波長・周期・振動数と $c=f\lambda$ を使いこなせる。 & \pageref{ph:波-5:sec2} & \bkcb & \bkcb \\
 & 正弦波の式 $y=A\sin(\omega t-kx)$ からゆれを求められる。 & \pageref{ph:波-5:sec3} & \bkcb & \bkcb \\
 & 媒質の運動方程式から波動方程式と波の速さを導ける。 & \pageref{ph:波-5:sec4} & \bkcb & \bkcb \\
 & おもりの列の分散関係・群速度と、計算の刻みの条件がわかる。 & \pageref{ph:波-5:sec7} & \bkcb & \bkcb \\
{\sffamily\bfseries\color{\bkTT}波-6} & 固定端・自由端の反射と、定在波の腹と節を説明できる。 & \pageref{ph:波-6:sec1} & \bkcb & \bkcb \\
 & 弦の固有振動数を求め、フレットと音の高さの関係がわかる。 & \pageref{ph:波-6:sec3} & \bkcb & \bkcb \\
 & 閉管・開管の固有振動数と共鳴から、音の速さを求められる。 & \pageref{ph:波-6:sec4} & \bkcb & \bkcb \\
 & 変数分離とフーリエ級数で、はじいた弦の倍音を求められる。 & \pageref{ph:波-6:sec5} & \bkcb & \bkcb \\
 & はりの振動の式を解き、鉄琴や配管の固有振動数を求められる。 & \pageref{ph:波-6:sec9} & \bkcb & \bkcb \\
{\sffamily\bfseries\color{\bkTT}波-7} & 波が運ぶエネルギーと強さを、振幅・振動数・距離から求められる。 & \pageref{ph:波-7:sec1} & \bkcb & \bkcb \\
 & インピーダンスの比から、境目ではね返る割合と通りぬける割合を求められる。 & \pageref{ph:波-7:sec4} & \bkcb & \bkcb \\
 & 表計算で、2本の弦をつないだときの山のはね返りと通りぬけを追える。 & \pageref{ph:波-7:sec4} & \bkcb & \bkcb \\
 & 波の式と境界条件から反射係数を導き、固定端・自由端・はね返らない端を1つの式で説明できる。 & \pageref{ph:波-7:sec5} & \bkcb & \bkcb \\
 & 吸収による弱まり、整合層、光の反射率を、インピーダンスと dB で見積もれる。 & \pageref{ph:波-7:sec7} & \bkcb & \bkcb \\
{\sffamily\bfseries\color{\bkTT}波-8} & 近い高さの2つの音が重なるときの、うなりの回数を求められる。 & \pageref{ph:波-8:sec2} & \bkcb & \bkcb \\
 & ホイヘンスの原理で、反射・屈折の法則と回折を説明できる。 & \pageref{ph:波-8:sec1} & \bkcb & \bkcb \\
 & 2つの波源からの道のりの差で、強め合う所と弱め合う所を求められる。 & \pageref{ph:波-8:sec3} & \bkcb & \bkcb \\
 & 波面を数えてドップラー効果の式を導き、聞こえる高さを求められる。 & \pageref{ph:波-8:sec5} & \bkcb & \bkcb \\
 & 斜めのドップラー効果・反射板・衝撃波で、速さをはかるしくみを説明できる。 & \pageref{ph:波-8:sec7} & \bkcb & \bkcb \\
{\sffamily\bfseries\color{\bkTT}波-9} & 雷やこだまの時間から、音の速さを使って距離を求められる。 & \pageref{ph:波-9:sec2} & \bkcb & \bkcb \\
 & 音の大きさを dB で表し、音を合わせたときや、離れたときの大きさを計算できる。 & \pageref{ph:波-9:sec7} & \bkcb & \bkcb \\
 & 縮みにくさと密度から気体・液体・固体の音の速さを求め、遮音と残響を見積もれる。 & \pageref{ph:波-9:sec3} & \bkcb & \bkcb \\
{\sffamily\bfseries\color{\bkTT}波-10} & こだまの時間から、コウモリ・ソナー・超音波探傷の距離を求められる。 & \pageref{ph:波-10:sec1} & \bkcb & \bkcb \\
 & 振動の変位・速度・加速度を換算し、回転機械のスペクトルから故障の見当をつけられる。 & \pageref{ph:波-10:sec6} & \bkcb & \bkcb \\
 & 1自由度の模型で、加速度計と変位計のしくみを説明し、感度のずれを計算できる。 & \pageref{ph:波-10:sec7} & \bkcb & \bkcb \\
{\sffamily\bfseries\color{\bkTT}波-11} & 音をサンプリングと量子化で数の列にし、データの量を求められる。 & \pageref{ph:波-11:sec1} & \bkcb & \bkcb \\
 & エイリアシングで、車輪が逆に回って見えるわけを説明できる。 & \pageref{ph:波-11:sec2} & \bkcb & \bkcb \\
 & 目もり $\Delta f=f_{\mathrm s}/N$ を読み、窓関数でもれを小さくできる。 & \pageref{ph:波-11:sec4} & \bkcb & \bkcb \\
 & サンプリング定理・フィルター・平均で、S/N 比を見積もれる。 & \pageref{ph:波-11:sec5} & \bkcb & \bkcb \\
 & 情報量と通信路容量を計算し、熱のエントロピーと比べられる。 & \pageref{ph:波-11:sec9} & \bkcb & \bkcb \\
{\sffamily\bfseries\color{\bkTT}波-12} & 光の直進・反射・屈折を、身近な例と作図で説明できる。 & \pageref{ph:波-12:sec1} & \bkcb & \bkcb \\
 & 屈折の法則で、屈折角・臨界角・見かけの深さを求められる。 & \pageref{ph:波-12:sec3} & \bkcb & \bkcb \\
 & レンズの式で、像の位置と倍率を求め、光学機器の倍率を見積もれる。 & \pageref{ph:波-12:sec6} & \bkcb & \bkcb \\
 & 屈折率が色で変わること（分散）から、プリズムと虹の色を説明できる。 & \pageref{ph:波-12:sec8} & \bkcb & \bkcb \\
 & フェルマーの原理から反射・屈折の法則を導き、光線の行列でレンズの組を計算できる。 & \pageref{ph:波-12:sec7} & \bkcb & \bkcb \\
{\sffamily\bfseries\color{\bkTT}波-13} & ヤングの実験・回折格子・薄い膜で、光が強め合う所と弱め合う所を、経路の差から計算できる。 & \pageref{ph:波-13:sec2} & \bkcb & \bkcb \\
 & 円い穴の回折で見分けられる角の限界と、偏光板を通る光の強さ・ブルースター角を計算できる。 & \pageref{ph:波-13:sec5} & \bkcb & \bkcb \\
 & ゆすられる電子のモデル（強制振動）から、空が青く夕焼けが赤いわけを説明できる。 & \pageref{ph:波-13:sec8} & \bkcb & \bkcb \\
{\sffamily\bfseries\color{\bkTT}波-14} & 電波の振動数からアンテナの大きさを、電波が往復する時間から距離を、見積もれる。 & \pageref{ph:波-14:sec5} & \bkcb & \bkcb \\
 & AM・FM・ディジタルの変調を、波の振幅・振動数・位相を変えることとして説明できる。 & \pageref{ph:波-14:sec2} & \bkcb & \bkcb \\
 & たくさんの周波数で同時に送る方式（OFDM）の速さを計算し、通信路容量の上限とくらべられる。 & \pageref{ph:波-14:ex3} & \bkcb & \bkcb \\
\end{longtblr}
\endgroup


%% ---------------- つぎの巻へ ----------------
\bkbacktitle{つぎの巻へ}
\newcommand\bknextpill[1]{\bk@pill{#1}}
{\small 波動で、ゆれと波を式で書き、共振・制御・音・光を調べる道具がそろいました。物理シリーズの『電磁気学』では、電気と磁気の法則を導きます。この巻で「同じ形の式」として使った回路や、光が電磁波であることが、そこで導かれます。\par}\vspace{3mm}
\begin{tcolorbox}[enhanced, colframe=phVol4, colback=white, boxrule=0.9pt, arc=2.5mm, left=3mm, right=3mm, top=2mm, bottom=2mm,
  title={物理シリーズの『電磁気学』：場から回路・機械・電波へ}, fonttitle=\sffamily\bfseries, colbacktitle=phVol4, coltitle=white]
\small
\begin{itemize}[leftmargin=1.2\zw, itemsep=2pt, topsep=1pt, label={\color{phVol4}●}]
\item \bk@pill{電-9}\ スイッチを入れた瞬間：過渡現象（波-1 ⑩の相似の表の RLC 回路を、回路の法則から導く）
\item \bk@pill{電-13}\ 長い電線を読む：分布定数回路（波-7 のインピーダンスと整合を、電線の上の波で）
\item \bk@pill{電-15}\ 光も電波も同じ波：電磁波（光が電磁波であること（波-12）と、波-13 で結果だけ使った「出す光の強さは加速度の2乗に比例」を導く）
\end{itemize}
\end{tcolorbox}
\phbkhead{この巻の道具は、この先でこう使う}
{\small この巻の章と、その道具を使う物理シリーズのほかの巻（『電磁気学』から先）の章です（章のはじめの「道しるべ」の「この章の先に」から自動で集めた表）。\par}\vspace{2mm}
% 自動生成（tools/mknext.py p3）。各章の道しるべ「この章の先に」から、後ろの巻の章を集めたもの（第13回）
\begin{longtblr}{colspec={Q[l,m,wd=58mm] X[l,m]}, width=\linewidth, rowhead=1, row{1}={bg=bkVol!14, font=\sffamily\bfseries\footnotesize},
  hlines={0.4pt,bkLine}, rowsep=1.5pt, cells={font=\footnotesize}}
この巻の章 & この道具を使う章 \\
{\bknextpill{波-1}\ ばねとおもりで考える：1自由度の振動} & {\bknextpill{電-9}\ スイッチを入れた瞬間：過渡現象} \\
{\bknextpill{波-7}\ 波が運ぶエネルギー：インピーダンスと反射} & {\bknextpill{電-13}\ 長い電線を読む：二端子対回路と分布定数回路} \\
{\bknextpill{波-12}\ 光の進み方：反射・屈折・レンズ} & {\bknextpill{電-15}\ 光も電波も同じ波：電磁波\\ \bknextpill{原-2}\ 光はつぶでもある：光電効果と光子} \\
{\bknextpill{波-13}\ 光の波：干渉・回折・偏光・散乱} & {\bknextpill{電-15}\ 光も電波も同じ波：電磁波\\ \bknextpill{原-2}\ 光はつぶでもある：光電効果と光子} \\
\end{longtblr}


%% ---------------- さくいん ----------------
\bkbacktitle{さくいん}
{\small 数字はページです。その語がはじめて出てくるページや、詳しく説明しているページをのせています。}\par\vspace{2mm}
\raggedcolumns\begin{multicols}{2}\small\setlength{\parskip}{0pt}
\directlua{bk.printidx()}
\end{multicols}

%% ---------------- 奥付 ----------------
%% さくいんのあとのメモのページの柱を「メモ」にする（第1・2巻と同じ）
\clearpage\bksetrunning{}{メモ}{メモ}%
\bkpadfour
\global\bk@decofalse\bkfolio{false}
\null\vfill
\noindent\begin{minipage}{\linewidth}\small
{\sffamily\bfseries\Large\color{bkDeep}\phSeries　第3巻}\par\vspace{1mm}
{\sffamily\bfseries\large 波動：ゆれと波を読む}\par\vspace{3mm}
{\color{bkVol}\hrule height 1pt}\vspace{3mm}
\begin{tabular}{@{}l@{\hspace{6mm}}l@{}}
編者 & \phEditor \\
 & 委員長　小学校長（まとめ役） \\
 & 委員　数学・電気工学・機械工学・教育学・量子力学・建築工学・\\
 & \phantom{委員　}半導体工学の大学教授、幼稚園教諭、塾講師、\\
 & \phantom{委員　}特別支援学校教諭、心理士、科学解説者、小学校教諭 \\
デザイン & ユニバーサルデザイン担当委員 \\
挿絵 & イラスト担当委員 \\
先に読む本 & 「物理を理解するための数学」全9巻、物理シリーズ第1〜2巻 \\
版 & 第1版 Ver0（作成中）　2026年11月 \\
\end{tabular}\par\vspace{3mm}
{\color{bkVol}\hrule height 1pt}\vspace{2mm}
{\footnotesize この本の問題の答えは、巻末の「こたえ」にあります。詳しい解き方は、別冊の問題集と指導書にのせる予定です。物性値の出典は、巻末の「音速・屈折率・物性値」の表にあります。}
\end{minipage}
\makeatother
